% Chapter 1, typeset from lectures/L01: the README is the opener and the closing review, and each
% part of Appendix A is a section, in the same order and with the same subsections. Appendix B,
% the exercises, is the last section.

\chapter{Combinational Logic and First VHDL}
\label{c1:ch}

{\large\itshape From a truth table to a working gate network, in CircuitVerse and in VHDL.\par}

\begin{chapterbox}{In this chapter}
\begin{itemize}
  \item Logic gates and truth tables, and why digital logic tolerates noise that analog cannot.
  \item Boolean algebra: reading a sum-of-products equation straight off a truth table.
  \item Building and simulating a gate network by hand in CircuitVerse.
  \item VHDL from the first chapter: \code{entity}, \code{architecture}, \code{std_logic}, and the
    concurrent assignment.
  \item A braking assistant, from a requirement and its truth table to a gate network, and on to
    switches and an LED on a DE0-CV FPGA board.
  \item \code{or_gate}, taken from truth table to complete module in \secref{c1:a4} and
    \secref{c1:a5}.
\end{itemize}
\end{chapterbox}

This chapter takes one circuit from a requirement to running hardware, and then does it again. The
first time is \code{or_gate}, the smallest complete VHDL design there is, which \secref{c1:a4} and
\secref{c1:a5} take end to end in writing. The second is a braking assistant: four gates, derived
from a requirement and its truth table, simulated against it in CircuitVerse, written as an
\code{entity} and an \code{architecture} with an internal \code{signal} and two concurrent
assignments, and finally taken through Quartus onto a DE0-CV, with its four inputs on switches and
its brake output on an LED. Flick the switches and watch the car decide to brake.

In the course this book comes from, the braking assistant is built live, and \secref{c1:a6} gives
you its specification without its answer so that you can derive it yourself first, the way the
people in the room are asked to. Do that before you read on: it is the habit the rest of the book
is built on. You then write its VHDL in Exercise~\ref{c1:ex:7}, and a testbench decides whether your
reading of the requirement matches the truth table.

You need nothing installed to read this chapter. The exercises need GHDL, which
\secref{c2:app:c} installs and explains how to run; the three commands in \secref{c1:a5} are
enough to get started. The board work is shown in Appendix~\ref{app:quartus} for anyone who has a
DE0-CV; no exercise needs one.

% Sections 1.1 to 1.6, from lectures/L01/appendix/a_combinational_logic.md: one section per part,
% A.1 to A.6.

\section{Digital Signals and Logic Gates}
\label{c1:a1}
A digital signal takes one of two values, \code{0} or \code{1}, rather than the continuous range an
analog signal uses. Anything close enough to \code{0} or \code{1} is read as exactly \code{0} or
\code{1}, so small amounts of electrical noise change nothing. That noise immunity is why almost all
modern computation, from microcontrollers to FPGAs, is digital.

A \textbf{logic gate} has one or more single-bit inputs, one single-bit output, and a fixed Boolean
function between them.

\textbf{Combinational logic}, this chapter's topic, produces outputs from the current inputs alone,
with no memory. \textbf{Sequential logic} (Chapter~\ref{c3:ch}) depends on the current inputs
\emph{and} on state stored in flip-flops.

Every basic 2-input gate, plus the unary \code{NOT}:

\begin{narrowtable}{@{}cccccccc@{}}
  \thead{A} & \thead{B} & \thead{AND} & \thead{OR} & \thead{NAND} & \thead{NOR} & \thead{XOR} &
    \thead{XNOR} \\
  \midrule
  0 & 0 & 0 & 0 & 1 & 1 & 0 & 1 \\
  0 & 1 & 0 & 1 & 1 & 0 & 1 & 0 \\
  1 & 0 & 0 & 1 & 1 & 0 & 1 & 0 \\
  1 & 1 & 1 & 1 & 0 & 0 & 0 & 1 \\
\end{narrowtable}

\begin{narrowtable}{@{}cc@{}}
  \thead{A} & \thead{NOT A} \\
  \midrule
  0 & 1 \\
  1 & 0 \\
\end{narrowtable}

\noindent Worth internalizing straight from the tables:
\begin{itemize}
  \item \code{NAND}, \code{NOR} and \code{XNOR} are the inverses of \code{AND}, \code{OR} and
    \code{XOR}.
  \item \code{XOR} is \code{1} when the inputs differ. Beyond two inputs this generalizes to
    \textbf{odd parity}: \code{1} when an odd number of inputs are \code{1}
    (Exercise~\ref{c1:ex:3}).
  \item Any Boolean function can be built from \code{AND}, \code{OR} and \code{NOT}, and also from
    \code{NAND} alone or \code{NOR} alone.
\end{itemize}

\section{Boolean Algebra Notation}
\label{c1:a2}
Three operators, mapping directly onto the gates above:

\begin{narrowtable}{@{}ll@{}}
  \thead{Gate} & \thead{Boolean algebra notation} \\
  \midrule
  AND & multiplication: \code{A} \textbf{AND} \code{B} is written \code{AB} \\
  OR & addition: \code{A} \textbf{OR} \code{B} is written \code{A + B} \\
  NOT & prime: \textbf{NOT} \code{A} is written \code{A'} \\
\end{narrowtable}

\noindent A function written as an OR of AND-terms, such as \code{X = AB' + CD}, is in
\textbf{sum-of-products} form.

Every truth table reads off into sum-of-products directly:
\begin{enumerate}
  \item Find every row where the output is \code{1}.
  \item For each, write one AND-term containing every input: directly if it is \code{1} in that
    row, primed if it is \code{0}.
  \item OR those terms together.
\end{enumerate}
For example:

\begin{narrowtable}{@{}ccc@{}}
  \thead{A} & \thead{B} & \thead{X} \\
  \midrule
  0 & 0 & 0 \\
  0 & 1 & 1 \\
  1 & 0 & 1 \\
  1 & 1 & 0 \\
\end{narrowtable}

\noindent Row \code{AB = 01} contributes \code{A'B} and row \code{AB = 10} contributes \code{AB'},
giving \code{X = A'B + AB'}. Check it against the table: that is \code{XOR}, so the same network is
one \code{XOR} gate rather than two \code{AND} gates, two \code{NOT} gates and an \code{OR}.

\subsection{The laws, as a reminder}
\label{c1:laws}
You have met these. They are here so the book has one notation for them and so the later chapters
can cite them, not because they need teaching:

\begin{narrowtable}{@{}lll@{}}
  \thead{Law} & \thead{With OR} & \thead{With AND} \\
  \midrule
  Identity & \code{A + 0 = A} & \code{A · 1 = A} \\
  Null & \code{A + 1 = 1} & \code{A · 0 = 0} \\
  Idempotent & \code{A + A = A} & \code{A · A = A} \\
  Complement & \code{A + A' = 1} & \code{A · A' = 0} \\
  Absorption & \code{A + AB = A} & \code{A(A + B) = A} \\
  De Morgan & \code{(A + B)' = A'B'} & \code{(AB)' = A' + B'} \\
\end{narrowtable}

\noindent\textbf{De Morgan is the one this book leans on}, less for minimization than as a rule
about moving an inversion through a gate: an \code{OR} with both inputs inverted \emph{is} a
\code{NAND}, and an \code{AND} with both inputs inverted \emph{is} a \code{NOR}. It is why a
\code{NAND} gate can build anything, which Exercise~\ref{c1:ex:6} has you write in VHDL, and it
comes back every time an active-low signal meets logic written for active-high. In a book where
every reset and every button is active-low, that is often.

That gap between a straightforward sum-of-products reading and the smallest network is what a
\textbf{Karnaugh map} closes, and Chapter~\ref{c2:ch} works one end to end. Not because the
technique is new to you, but because the minimized equation is what gets written in VHDL, and
Chapter~\ref{c8:ch} derives its state machine's next-state logic exactly this way.

\section{Building and Simulating a Circuit in CircuitVerse}
\label{c1:a3}
CircuitVerse (\url{https://circuitverse.org/simulator}) is a free, browser-based logic simulator.
Every gate network in this book is built and tested there by hand before it is written in VHDL:
\begin{enumerate}
  \item Open the simulator and start a new project.
  \item Place an \textbf{input} element per input signal and an \textbf{output} element per output,
    renamed to match your variables.
  \item Place the gates your equation calls for, from the \emph{Gates} panel.
  \item Wire it up one gate at a time, following your derived equation. You should already know
    exactly which gates you need before you open the simulator, which is why \secref{c1:a2} comes
    first.
  \item Toggle each input and confirm the output matches your truth table for \textbf{every} row.
  \item Save via \code{Project -> Save Online}, or \code{Project -> Download as image/data} for a
    re-importable \file{.cv} file.
\end{enumerate}
\textbf{Step 5 is not optional, and nothing else does it for you.} Three rules make it a check
rather than a formality:
\begin{itemize}
  \item \textbf{Exhaustive, not spot-checked.} Two inputs is four rows, three is eight. At these
    sizes there is no excuse for sampling. For the sequential circuits later on this becomes ``step
    the clock and check at every edge'', the same discipline applied to time.
  \item \textbf{Derive first, simulate second.} Draw the circuit and \emph{then} decide what it
    should do, and you will read your expectations off the simulation, which proves nothing.
  \item \textbf{A mismatch is information.} It says the drawing and the equation disagree, and one
    of them is wrong.
\end{itemize}
Every hardware bug you will chase is found by predicting a behaviour and then checking it. Doing
it by hand on a four-gate circuit is how you learn to do it on one you cannot see all at once.

Once the design is VHDL, a self-checking testbench takes over that job. Every worked example and
every exercise ships one; you run them rather than write them, as covered in \secref{c2:app:c}.

\section{Enough VHDL to Read and Write a Gate Circuit}
\label{c1:a4}
VHDL describes hardware, not a sequence of instructions. Every statement below runs concurrently
and continuously, exactly like the gates it describes. The language arrives as it is needed; this
section is enough for a simple combinational circuit.

Every VHDL module has two parts:
\begin{itemize}
  \item An \textbf{entity}, the black-box view: the ports, and nothing about what happens inside.
  \item An \textbf{architecture}, the implementation behind it.
\end{itemize}
That split is a contract, not ceremony. The entity is everything another design needs in order to
use this module, which is what lets one design instantiate another without reading it, and lets
you rewrite an implementation without touching what depends on it. The shape is a header beside a
source file, or an interface beside its class; the difference is that VHDL enforces it for every
module, with no way to skip the interface or leak the implementation through it.

Take \code{or_gate}: inputs \code{a}, \code{b}, output \code{x}, and \code{x = a or b}.

\begin{vhdlcode}
entity or_gate is
    port(a, b: in std_logic;
         x   : out std_logic);
end entity;
\end{vhdlcode}

\bookfigure{or_gate_entity}{The entity \code{or_gate}: two inputs and an output, and nothing about
what happens inside.}{c1:fig:entity}

\noindent\code{std_logic} is the single-bit signal type. It is not built into VHDL but comes from
the \code{std_logic_1164} package, which is why every file in this book opens with:

\begin{vhdlcode}
library ieee;
use ieee.std_logic_1164.all;
\end{vhdlcode}

\noindent It has nine values rather than two, because real hardware needs more than true and false:
\begin{itemize}
  \item \code{'0'} and \code{'1'}: driven low and driven high. The only two you will write in this
    book.
  \item \code{'Z'}: high-impedance, nothing driving the line at all.
  \item \code{'U'}: uninitialized, what a signal holds before anything assigns it.
  \item \code{'X'}: unknown, what a simulator shows when two things drive one signal and disagree.
  \item \code{'-'}, \code{'W'}, \code{'L'}, \code{'H'}: don't-care and weak-drive values, rarely
    written by hand.
\end{itemize}
Recognizing the seven you will not write still matters: they are how a simulator tells you
something is wrong. A signal sitting at \code{'U'} or \code{'X'} is a bug, not a value.

The architecture fills in the box:

\begin{vhdlcode}
architecture behaviour of or_gate is
begin
    x <= a or b;
end architecture;
\end{vhdlcode}

\bookfigure{or_gate_arch}{The architecture of \code{or_gate}: the gate that fills the
box.}{c1:fig:arch}

\noindent Together they form the complete module:

\bookfigure{or_gate_module}{The module \code{or_gate}: the entity and the architecture
together.}{c1:fig:module}

\noindent\code{x <= a or b;} is \secref{c1:a2}'s Boolean algebra almost verbatim: VHDL's
\code{and}, \code{or}, \code{not} and \code{xor} work directly on \code{std_logic}, so any equation
you derived by hand becomes VHDL with barely any translation.

The assignment is \textbf{concurrent}. It is not executed once; it holds continuously, like a wire
between real gates.

The keyword \code{signal} declares an internal wire inside an architecture, as opposed to an
externally visible port. You use one in Exercise~\ref{c1:ex:7} to name an intermediate result
rather than nesting it inline, and Chapter~\ref{c2:ch} leans on the same idea throughout.

\section{The Complete Module: \code{or_gate}}
\label{c1:a5}
\Secref{c1:a4}'s pieces together make a complete, synthesizable file. This is the whole of
\file{lectures/L01/or_gate/or_gate.vhd}, and it is a real design that synthesizes and runs on the
DE0-CV:

\vhdlfile{lectures/L01/or_gate/or_gate.vhd}

\begin{itemize}
  \item The \code{library}/\code{use} clauses pull in \code{std_logic}.
  \item \code{entity} declares the outside: two inputs, one output.
  \item \code{architecture} implements the inside, as one concurrent assignment.
\end{itemize}
That is the smallest complete VHDL design there is, and the shape never changes. Every module in
this book, up to the state machines in Chapter~\ref{c8:ch}, is this same
\code{library}/\code{entity}/\code{architecture} skeleton. What grows is the architecture body,
never the frame.

\textbf{Checking your work.} \code{or_gate} ships \file{or_gate_tb.vhd}, which drives all four
input combinations and checks the output:

\begin{shell}
cd lectures/L01/or_gate
ghdl -a --std=93 or_gate.vhd or_gate_tb.vhd
ghdl -e --std=93 or_gate_tb
ghdl -r --std=93 or_gate_tb --assert-level=error --stop-time=10ms
\end{shell}

\noindent No output beyond a final note means every check passed. Testbenches are covered properly
in \secref{c2:app:c}; for now, treat these three commands as how you confirm a module works.

\section{The Worked Circuit: A Braking Assistant}
\label{c1:a6}
The course builds a different circuit from \code{or_gate}, live, so you see the same path walked
twice: once here in writing, and once on something you have not already read the answer to.

The requirement: \textbf{the car brakes if the driver brakes, or if the assistance system decides
to.} The assistance system decides to when either detector sees something, unless it is reporting a
fault, in which case it is not trusted to brake at all.

\begin{booktable}{@{}p{7.5em}p{5em}L@{}}
  \thead{Port} & \thead{Direction} & \thead{Meaning} \\
  \midrule
  \code{driver_brake} & in & The driver is on the brake pedal. \\
  \code{sensor} & in & The proximity sensor sees an obstacle. \\
  \code{radar} & in & The radar sees an obstacle. \\
  \code{error} & in & The assistance system is reporting a fault. \\
  \code{engine_brake} & out & Apply the brakes. \\
\end{booktable}

\noindent Three clauses decide the whole structure between them:
\begin{itemize}
  \item either detector is enough on its own, and neither is trusted more than the other.
  \item a fault does not raise an alarm; it \emph{suppresses} the assistance system's ability to
    brake.
  \item the driver can always brake, including when the assistance system is faulty and doing
    nothing.
\end{itemize}
That last one is a safety requirement, and it is the one to hold on to when you derive the
network: the pedal is not an input to the assistance logic, it bypasses it. Work out what changes
if the pedal is fed through the fault logic instead, and you have found the bug the structure
exists to prevent.

\subsection{The truth table}
\label{c1:adastable}
Sixteen rows, in the order a four-bit counter would produce them. \code{adas_brake} is not a port;
it is the assistance system's own decision, shown as an intermediate column so you can see where
the structure comes from.

\begin{narrowtable}{@{}cccccc@{}}
  \thead{\code{driver_brake}} & \thead{\code{sensor}} & \thead{\code{radar}} & \thead{\code{error}}
    & \thead{\code{adas_brake}} & \thead{\code{engine_brake}} \\
  \midrule
  0 & 0 & 0 & 0 & 0 & 0 \\
  0 & 0 & 0 & 1 & 0 & 0 \\
  0 & 0 & 1 & 0 & 1 & 1 \\
  0 & 0 & 1 & 1 & 0 & 0 \\
  0 & 1 & 0 & 0 & 1 & 1 \\
  0 & 1 & 0 & 1 & 0 & 0 \\
  0 & 1 & 1 & 0 & 1 & 1 \\
  0 & 1 & 1 & 1 & 0 & 0 \\
  1 & 0 & 0 & 0 & 0 & 1 \\
  1 & 0 & 0 & 1 & 0 & 1 \\
  1 & 0 & 1 & 0 & 1 & 1 \\
  1 & 0 & 1 & 1 & 0 & 1 \\
  1 & 1 & 0 & 0 & 1 & 1 \\
  1 & 1 & 0 & 1 & 0 & 1 \\
  1 & 1 & 1 & 0 & 1 & 1 \\
  1 & 1 & 1 & 1 & 0 & 1 \\
\end{narrowtable}

\noindent The bottom half is entirely \code{1}: once \code{driver_brake} is set nothing else changes
the answer. That is what ``overrides everything'' looks like in a table.

\subsection{Before you read on}
\label{c1:adasbefore}
Everything above is the specification. The derivation, the gate network, and the VHDL are what the
course builds live, and it is worth arriving at them with your own answer to compare against rather
than one you have already read. So, from the table above:
\begin{itemize}
  \item Read a sum-of-products equation for \code{engine_brake} straight off it, the way
    \secref{c1:a2} did.
  \item Simplify it by hand. The requirement is three clauses long, so the minimal network is a
    great deal smaller than the row-by-row reading gives you. Work out how much smaller before you
    are told.
  \item Count the gates your simplified equation needs, and write the number down.
  \item Sketch the network, and decide which single gate the safety argument above rests on.
\end{itemize}
Write those four answers down. Being wrong about any of them is worth more than not having guessed,
because you will know exactly which step to re-examine when you check them. That is
\secref{c1:a3}'s ``derive first, simulate second'' applied to a circuit you have not built yet, and
it is the habit the rest of the book is built on. Afterwards you write the VHDL yourself, in
Exercise~\ref{c1:ex:7}, and a testbench decides whether your equations and this truth table
agree.

\subsection{Why there is a testbench for it}
\label{c1:adastb}
Sixteen combinations is few enough to check by hand once. It is not few enough to check every time
somebody edits the file, which is the situation every module in this book is in from here on.

A \textbf{testbench} is a second VHDL file that does the checking. It is not part of the design and
never reaches the FPGA: it drives the inputs through every case, works out what the output should
have been, compares, and reports. For \code{adas}, the braking assistant's module, that is a loop
over the sixteen combinations restating the braking rule independently.

\emph{Independently} is the important word. A testbench that read its expected answer out of the
design would agree with it by construction, including when the design is wrong. Restating the
requirement separately is what turns it into a check.

That is why almost every directory in the companion repository has a \file{_tb.vhd} beside the
module: it is how you can be handed a module, change it, and know within a second whether you broke
it. \Secref{c2:app:c} covers running them; you are not asked to write one anywhere in this book.

% Section 1.7, from lectures/L01/README.md: the objectives, the questions to test yourself, and the
% next lecture.

\needspace{10\baselineskip}
\section{Review}
\label{c1:sec:review}

\subsection*{What you should be able to do now}
After this chapter, you should be able to:
\begin{itemize}
  \item Derive a sum-of-products equation from a truth table, then build and simulate the network in
    CircuitVerse.
  \item Build \code{NOT}, \code{AND} and \code{OR} from \code{NAND} gates alone, and use De Morgan
    to explain why the \code{OR} case works.
  \item Read and write a combinational VHDL module, an \code{entity} with \code{std_logic} ports and
    an \code{architecture} driving its output, and say why VHDL separates the two.
  \item Explain why a concurrent assignment describes a wire rather than a statement that runs once,
    and why a signal sitting at \code{'U'} or \code{'X'} is a bug rather than a value.
  \item Run a provided testbench and tell whether your module passed.
\end{itemize}

\needspace{6\baselineskip}
\subsection*{Questions to test yourself}
You should be able to explain:
\begin{enumerate}
  \item What is the difference between an \code{entity} and an \code{architecture}, and why does
    VHDL separate them?
  \item Why is \code{x <= a or b;} not ``a statement that runs once''? What does it describe
    instead?
  \item In the braking assistant, the driver's pedal is ORed in last rather than passing through
    the fault logic. What breaks if you wire it the other way?
\end{enumerate}

\needspace{6\baselineskip}
\subsection*{Looking ahead}
Chapter~\ref{c2:ch} is about larger networks:
\begin{itemize}
  \item Karnaugh maps: finding the \emph{smallest} gate network, not just a correct one.
  \item \code{std_logic_vector}, \code{process} and \code{case}.
  \item Submodules, demonstrated on a two-digit hexadecimal display driven by eight switches.
  \item Self-checking testbenches: verifying your own VHDL without a board.
\end{itemize}

% Section 1.8, from lectures/L01/appendix/b_exercises.md.

\section{Exercises}
\label{c1:sec:exercises}

\begin{aside}
\textbf{How to work these.} Exercise~\ref{c1:ex:1} is a short pen-and-paper warm-up on notation this
book uses throughout; it needs no VHDL, GHDL, or board. Each exercise is tagged with the kinds of
work it asks for: \kinds{Paper} by hand, \kinds{CircuitVerse} in the browser, \kinds{VHDL} in a
module of your own.

\textbf{How to check your work.} From Exercise~\ref{c1:ex:2} on, every exercise that asks you to
write a VHDL module ships a self-checking testbench under \file{lectures/L01/exercises} in the
companion repository. Write your module in its \file{exercises/<module>} directory, using the entity
name and the \textbf{port order} the exercise specifies, then run it with GHDL; \secref{c2:app:c}
has the three commands.

No FPGA board is needed for any exercise. Your job is to get the VHDL right, and the testbench is
how you confirm it. Short answers to the parts you work on paper are in
Appendix~\ref{app:answers}; reference solutions to the VHDL are published in the companion
repository, under \file{lectures/L01/solutions}. Try each exercise yourself before looking at
either.
\end{aside}

% ------------------------------------------------------------------------------------------
\exerciseset{Truth Tables and Gates}
\label{c1:set:1}

\exercise{A function that ignores an input}{Paper, CircuitVerse}
\label{c1:ex:1}
A function of two inputs \code{A}, \code{B} and one output \code{X} has the truth table below:

\begin{narrowtable}{@{}ccc@{}}
  \thead{A} & \thead{B} & \thead{X} \\
  \midrule
  0 & 0 & 1 \\
  0 & 1 & 0 \\
  1 & 0 & 1 \\
  1 & 1 & 0 \\
\end{narrowtable}

\begin{parts}
  \item Derive the sum-of-products equation for \code{X} directly from the truth table
    (\secref{c1:a2}). Two rows have \code{X = 1}, so you get two AND-terms.
  \item Simplify it algebraically. Factor out what the two terms share, and apply
    \code{A + A' = 1} and \code{1 · B' = B'} from \secref{c1:a2}. You should end up with a single
    literal.
  \item Your answer to \textbf{b)} says something about \code{A}. What? Which single gate from
    \secref{c1:a1}'s table computes the function, and what happens to the \code{A} input when you
    draw it?
  \item Build that single-gate version in CircuitVerse and confirm it reproduces every row of the
    truth table above. Leave the \code{A} input unconnected, or simply leave it out of the drawing;
    \textbf{b)} is the justification for doing so.
\end{parts}

% ------------------------------------------------------------------------------------------
\exerciseset{Entities and Architectures}
\label{c1:set:2}

\exercise{AND and NAND}{VHDL}
\label{c1:ex:2}
Write two small gates as complete VHDL modules, \code{and_gate} and \code{nand_gate}.

Neither gate is the point, and you should not expect either to be difficult. The point is that
this is where an entity and an architecture stop being something you read in \secref{c1:a4} and
become something you write, and that the two halves are worth writing in that order.

Both entities have the same ports:

\begin{narrowtable}{@{}lll@{}}
  \thead{Port} & \thead{Direction} & \thead{Type} \\
  \midrule
  \code{a} & in & \code{std_logic} \\
  \code{b} & in & \code{std_logic} \\
  \code{x} & out & \code{std_logic} \\
\end{narrowtable}

\begin{parts}
  \item Write the \code{entity} declaration for \code{and_gate}, and nothing else. In VHDL a
    module's interface is a thing you can write down and reason about before any implementation
    exists at all, which is the habit this part exists to build. Identify the required ports, the
    direction of each, and the type of each.
  \item Add an \code{architecture} implementing the AND gate with a single concurrent signal
    assignment.
  \item Now write \code{nand_gate}: the same interface with an inverted output, using VHDL's
    \code{nand} operator directly. Do \textbf{not} write \code{not (a and b)}.
\end{parts}
The two are equivalent as Boolean algebra, and on an FPGA they make no difference at all, because
both end up in the same lookup table. It matters on the kind of hardware Exercise~\ref{c1:ex:6} is
about, where NAND is the gate you physically have and everything else has to be built from it.
Writing what you mean, and letting the tools decide what it costs, is the habit; reaching past an
operator the language already gives you is the thing to avoid.

\modulefigure{and_gate}
\modulefigure{nand_gate}

\selfcheck name your entities \code{and_gate} and \code{nand_gate}, each with ports \code{a},
\code{b} (in) and \code{x} (out), declared in that order; their testbenches are in
\file{exercises/and_gate} and \file{exercises/nand_gate}, and each checks all four rows of the
corresponding truth table.

\exercise{Three-input XOR}{Paper, VHDL, CircuitVerse}
\label{c1:ex:3}
An XOR gate has three inputs, \code{a}, \code{b}, \code{c}, and one output, \code{x}. The output is
\code{1} whenever an odd number of the inputs are \code{1} (the 3-input XOR, or parity, function).

\begin{parts}
  \item Write the truth table for \code{x} over all eight combinations of \code{a}, \code{b},
    \code{c}, and confirm it matches the ``odd number of \code{1}s'' description.
  \item Implement the gate in VHDL, as a module of your own with an entity whose ports are three
    \code{std_logic} inputs and one \code{std_logic} output, and an architecture that drives
    \code{x} with a single concurrent assignment (\secref{c1:a4} shows the entity/architecture
    pattern and notes that VHDL's \code{xor} operator works directly on \code{std_logic};
    \secref{c1:a5} shows a complete module end to end).
  \item Build the same gate in CircuitVerse and confirm it reproduces every row of your truth
    table.
\end{parts}
\note[Tip] \code{xor} chains left-to-right, so the whole architecture body is a single line,
\code{x <= a xor b xor c;}, with no internal \code{signal} or intermediate gates needed.

\modulefigure{xor3}

\selfcheck name your VHDL entity \code{xor3}, with ports \code{a}, \code{b}, \code{c} (in) and
\code{x} (out), declared in that order; its testbench is in \file{exercises/xor3}, and
\secref{c1:a5} lists the three commands that run it.

\exercise{Majority gate}{Paper, VHDL, CircuitVerse}
\label{c1:ex:4}
A \textbf{majority gate} has three inputs, \code{a}, \code{b}, \code{c}, and one output, \code{x}.
The output is \code{1} whenever \emph{at least two} of the three inputs are \code{1}.

This is the circuit behind a triple-redundant vote: three sensors report the same measurement, and
the system acts on whatever answer two of them agree on, so a single failed sensor cannot decide
anything on its own.

\begin{parts}
  \item Write the truth table for \code{x} over all eight combinations of \code{a}, \code{b},
    \code{c}.
  \item Derive the sum-of-products equation for \code{x} directly from the truth table
    (\secref{c1:a2}):
    \begin{itemize}
      \item Four rows have \code{x = 1}, so you get four AND-terms.
      \item Each term includes all three inputs, primed where that input is \code{0} in its row.
    \end{itemize}
  \item Implement the gate in VHDL, as a module of your own:
    \begin{itemize}
      \item An entity with three \code{std_logic} inputs and one \code{std_logic} output.
      \item An architecture driving \code{x} with a single concurrent assignment, translating your
        four AND-terms directly (\secref{c1:a4}).
    \end{itemize}
  \item Build the same network in CircuitVerse and confirm it reproduces every row of your truth
    table.
\end{parts}
\note[Tip] Write the equation out on paper before you open either the editor or CircuitVerse. Four
three-input AND-terms plus one four-input OR gate is a lot of wiring to correct after the fact.

\note[Looking ahead] count the gates your equation needs, then convince yourself by inspection that
\code{x = ab + ac + bc} computes the same function with three two-input AND gates instead of four
three-input ones. Sum-of-products gave you a \emph{correct} network, not the \emph{smallest} one;
finding the smaller one systematically, rather than by staring at it, is what Chapter~\ref{c2:ch}
is for.

\modulefigure{majority3}

\selfcheck name your VHDL entity \code{majority3}, with ports \code{a}, \code{b}, \code{c} (in) and
\code{x} (out), declared in that order; its testbench is in \file{exercises/majority3} and checks
all eight rows. Either equation passes it: the testbench checks the function, not which network
you chose.

\exercise{Half adder}{VHDL, CircuitVerse}
\label{c1:ex:5}
Write a \textbf{half adder}: \code{sum} and \code{carry} from two single bits. You know the circuit;
the reason it is here is that it is the first module in the book with \textbf{two outputs}, and that
is a VHDL question rather than a logic one.

\begin{parts}
  \item Implement it, with two concurrent assignments and no internal \code{signal}.

    Two outputs does not mean two entities, two architectures, or a process. Each output port is
    driven by its own concurrent assignment, and the two are not steps that happen in an order:
    they describe two separate pieces of wire, both live all the time. If you find yourself
    reaching for a process to ``return'' two values, that instinct is the software one, and
    \secref{c1:a4}'s ``a concurrent assignment describes a wire'' is the correction.
  \item Build it in CircuitVerse and confirm both outputs, so you have seen the two-output shape as
    a drawing before you meet it as an entity with several \code{out} ports in
    Chapter~\ref{c2:ch}.
\end{parts}
\note[Looking ahead] this is the cell Chapter~\ref{c6:ch}'s counter is built from. The 4-bit adder
there is four of these chained carry-to-carry, and the counter is that adder with its output fed
back to its input, which is why the wraparound needs no logic of its own.

\modulefigure{half_adder}

\selfcheck name your entity \code{half_adder}, with ports \code{a}, \code{b} (in) and \code{sum},
\code{carry} (out), declared in that order; its testbench is in \file{exercises/half_adder} and
checks \code{sum} and \code{carry} separately, so a design that has one right and the other wrong is
told which is which.

% ------------------------------------------------------------------------------------------
\exerciseset{Universal Gates}
\label{c1:set:3}

\exercise{OR from NAND}{Paper, CircuitVerse, VHDL}
\label{c1:ex:6}
\Secref{c1:a1} states that any Boolean function can be built using only \code{NAND} gates. This
exercise makes you prove it to yourself for the three operators you have been using.

\begin{parts}
  \item Derive the \code{OR} case on paper before you build anything, because De Morgan
    (\secref{c1:laws}) gives it to you directly. Start from \code{(A'B')'} and apply the law; you
    should arrive at \code{A + B}. Now read that expression back as a circuit: what is \code{A'B'}
    with an inversion on its output, and how many \code{NAND} gates is the whole thing?
  \item In CircuitVerse, build each of the following from \code{NAND} gates and nothing else,
    confirming each against its truth table in \secref{c1:a1} as you go:
    \begin{itemize}
      \item \code{NOT a}: one NAND gate is enough, if you feed the same signal to both of its
        inputs.
      \item \code{a AND b}: one NAND gate, followed by the \code{NOT} you just built.
      \item \code{a OR b}: the network you derived in \textbf{a)}. Check your gate count against
        what you predicted.
    \end{itemize}
  \item Write the \code{OR} version in VHDL, as a module of your own:
    \begin{itemize}
      \item An entity with two \code{std_logic} inputs and one \code{std_logic} output.
      \item An architecture driving \code{x} with a single concurrent assignment.
      \item Use only VHDL's \code{nand} operator:
      \begin{itemize}
        \item no \code{or}, no \code{and}, no \code{not} anywhere in the expression.
      \end{itemize}
    \end{itemize}
  \item Explain why this matters in practice. Consider that a chip fabrication process has to
    physically implement every gate type it offers, and that \code{NAND} is the cheapest gate to
    build in CMOS.
\end{parts}
\note[Tip] VHDL's \code{nand} operator does not chain the way \code{and} and \code{or} do:
\code{a nand b nand c} is not legal. Parenthesize every NAND explicitly, which is what you want
here anyway, since each pair of parentheses is one physical gate in your CircuitVerse drawing.

\modulefigure{or_from_nand}

\selfcheck name your VHDL entity \code{or_from_nand}, with ports \code{a}, \code{b} (in) and
\code{x} (out), declared in that order; its testbench is in \file{exercises/or_from_nand} and checks
it against all four rows of the \code{OR} truth table.

% ------------------------------------------------------------------------------------------
\exerciseset{Putting It Together}
\label{c1:set:4}

\exercise{The braking assistant}{VHDL, Paper}
\label{c1:ex:7}
Write the braking assistant from \secref{c1:a6} in VHDL.

\Secref{c1:a6} had you derive its equations, count its gates and sketch its network \emph{before}
reading on. This is where you write it yourself, and where a testbench decides whether your reading
of the requirement matches the one in \secref{c1:a6}'s truth table.

The entity must have:

\begin{booktable}{@{}p{7.5em}p{4.8em}p{5.2em}L@{}}
  \thead{Port} & \thead{Direction} & \thead{Type} & \thead{Description} \\
  \midrule
  \code{driver_brake} & in & \code{std_logic} & The driver is on the brake pedal. \\
  \code{sensor} & in & \code{std_logic} & The proximity sensor sees an obstacle. \\
  \code{radar} & in & \code{std_logic} & The radar sees an obstacle. \\
  \code{error} & in & \code{std_logic} & The assistance system is reporting a fault. \\
  \code{engine_brake} & out & \code{std_logic} & Apply the brakes. \\
\end{booktable}

\begin{parts}
  \item Implement it from \textbf{your own} equations rather than from anyone else's. If your
    equations and the truth table disagree, that disagreement is the most useful thing this
    exercise can give you, and the testbench will tell you which way.

    Use two concurrent assignments and one internal \code{signal}:
    \begin{itemize}
      \item The signal holds the assistance system's own decision, before the driver's pedal is
        considered.
      \item This is the first exercise in the book with an internal signal, and it is the reason the
        keyword exists: a name for an intermediate result, so each line still reads like the clause
        of the requirement it came from.
    \end{itemize}
  \item Now break it deliberately, in the specific way \secref{c1:a6} warns about. Feed
    \code{driver_brake} \emph{through} the fault logic instead of letting it bypass:

\begin{vhdlcode}
engine_brake <= (driver_brake or sensor or radar) and (not error);
\end{vhdlcode}

    That is one gate simpler, and it looks reasonable. Run the testbench against it and answer:
    \begin{itemize}
      \item Which input combination does it fail on first, and what is physically happening in the
        car at that moment?
      \item How many of the sixteen rows does it get wrong, and what do they all have in common?
      \item The failing rows are exactly the ones the safety requirement was written for. What does
        that tell you about the value of a truth table you derived from a requirement, as opposed
        to one you read off a circuit you had already built?
    \end{itemize}
\end{parts}

\modulefigure{adas}

\selfcheck name your entity \code{adas}, with inputs \code{driver_brake}, \code{sensor},
\code{radar}, \code{error} and output \code{engine_brake}, declared in that order; its testbench is
in \file{exercises/adas} and checks all sixteen rows of \secref{c1:a6}'s truth table.

